% Clausura focal de antecedentes: fuentes del artículo XI, edición 21-09-2026.
% Rangos, copias íntegras y huellas: gestion/mapa_restitucion_memoria.json.
\section{Antecedentes de restitución operatoria y memoria}
\label{md:antmem:capitulo}

El corte de estos antecedentes está después de la producción conjunta
APP--TRIT--TPK del estado enriquecido y de la estructura discreta del
continuo. Se trabaja sobre el registro dodecafásico ya generado y sobre
realizaciones declaradas de sus lectores. El portador lineal posterior
y los candidatos de una reconstrucción no reemplazan el dominio de
historias generadas. Tampoco una realización reversible posterior se
identifica, por su sola notación, con la holonomía nonádica \(\Gamma_9\).

En la coordenatización y orientación del registro,
\begin{equation}
 \begin{gathered}
 K=(234,543,140,729,659,824,621,58,914,794,146,601),\\
 \mathbf1^{*}K=6263,\qquad
 (Sx)_i=x_{i+1},\quad i\pmod {12}.
 \end{gathered}
 \label{md:antmem:registro}
\end{equation}
El registro se usa después de su generación, no como entrada
retrospectiva para seleccionar una salida. El operador \(S\) es aquí
el desplazamiento posicional marcado. Sus dominios se mantendrán
separados de los planos canónicos utilizados más adelante.

\subsection{La órbita completa como referencia de su portador}
\label{md:antmem:k:orbita}

En \(V=\mathbb R^{12}\), con su producto euclídeo canónico, consideremos
\[
 O_K=(K,SK,\ldots,S^{11}K):\mathbb R^{12}\longrightarrow V.
\]
El desplazamiento \(S\) actúa sobre posiciones del registro. No se
identifica por la sola periodicidad con la evolución completa TPK o con
una simetría excepcional posterior.

\begin{theorem}[Referencia cíclica completa y Gram exacto]
\label{md:antmem:k:marco-ciclico}
La matriz \(O_K\) es invertible. Su Gram tiene los valores propios
\[
\begin{array}{c|c}
 \text{valor}&\text{multiplicidad}\\\hline
 39225169=6263^2&1\\
 1248025-345420\sqrt3&2\\
 589609&2\\
 1407529&2\\
 1053157&2\\
 1248025+345420\sqrt3&2\\
 697225&1
\end{array}
\]
y satisface
\begin{equation}
 589609\|c\|^2\leq\|O_Kc\|^2\leq39225169\|c\|^2,
 \quad
 \|O_K^{-1}\|=\frac1{\sqrt{589609}},\quad
 \operatorname{cond}(O_K)=\frac{6263}{\sqrt{589609}}.
 \label{md:antmem:k:cotas-orbita}
\end{equation}
La órbita del registro centrado tiene rango once.
\end{theorem}
\begin{proof}
Las entradas del Gram son las autocorrelaciones
\(\langle K,S^jK\rangle\). En orden \(j=0,\ldots,11\) valen
\begin{align*}
 (&4251257,2999323,3163385,3287920,3216553,3344743,\\
  &2950064,3344743,3216553,3287920,3163385,2999323).
\end{align*}
Como \(S\) es ortogonal, el Gram es circulante. Sustituir los doce
caracteres \(\exp(2\pi i m/12)\) en esa fila da la tabla real
anterior. Todos los valores son positivos. El menor de los dos valores
con \(\sqrt3\) es mayor que \(589609\), pues
\(658416>345420\sqrt3\), desigualdad que se comprueba elevando al
cuadrado sus dos miembros positivos. Los demás valores de la tabla son
mayores o iguales a \(589609\). El modo uniforme tiene valor
\((\sum K_i)^2=6263^2\); la desigualdad triangular de la transformada
de Fourier, usando \(K_i\geq0\), acota todos los demás por ese mismo
valor. El teorema espectral da \eqref{md:antmem:k:cotas-orbita}. Centrar
el registro elimina sólo el modo uniforme; los once restantes conservan
los valores no nulos de la tabla.
\end{proof}

El determinante, con el orden de columnas y la orientación de \(S\)
fijados aquí, es
\[
 \det O_K=5483119259214020183850397930008283625.
\]
La ausencia de un período propio del vector, por sí sola, no probaría la
invertibilidad: lo decisivo es que no se anula ninguno de sus modos de
Fourier. La condición numérica es aproximadamente \(8.15643\). Si se
normaliza \(K\) a norma uno, los extremos del Gram se dividen por
\(\|K\|^2=4251257\); la amplitud de sus componentes no debe
confundirse con una mejora intrínseca frente a toda referencia de igual
norma.

\begin{corollary}[Separación polar de forma y orientación]
\label{md:antmem:k:polar}
Sean
\[
 G_K=O_K^*O_K,\qquad P_K=G_K^{1/2},\qquad
 F_K=O_KG_K^{-1/2}.
\]
Entonces \(P_K\) es positivo definido, \(F_K\) es ortogonal y
\(O_K=F_KP_K\). Una realización isométrica \(J:V\to\mathcal Y\)
conserva \(G_K\) y \(P_K\), y transporta la parte isométrica a
\(JF_K\).
\end{corollary}
\begin{proof}
La positividad del Gram permite definir su raíz e inversa mediante el
teorema espectral. La identidad
\(F_K^*F_K=G_K^{-1/2}G_KG_K^{-1/2}=I\) prueba la afirmación.
Además, \((JO_K)^*(JO_K)=O_K^*O_K\).
\end{proof}

Esta polar se define a partir del Gram completo: no se obtiene sólo de la
lista de sus valores propios ni de la norma de \(K\). El símbolo
\(P_K\) de este corolario no es el proyector dimensional de rango diez
de otra realización incidencial. Esta separación es una consecuencia de la
referencia cíclica, no una identificación de todos sus lectores.

\subsection{Identificación de operadores y holonomía relativa}
\label{md:antmem:k:identificacion}

\begin{theorem}[Doce respuestas generales o una respuesta cíclica]
\label{md:antmem:k:operadores}
Para cualquier operador lineal \(H:V\to V\), sus doce respuestas
\(Y=(HK,HSK,\ldots,HS^{11}K)\) determinan
\[
 H=YO_K^{-1}.
\]
Si \(HS=SH\), una única respuesta vectorial completa \(y=HK\)
determina
\begin{equation}
 H=O_yO_K^{-1},\qquad
 H=\sum_{j=0}^{11}h_jS^j,\qquad h=O_K^{-1}y.
 \label{md:antmem:k:filtro}
\end{equation}
Con errores \(E\) en la matriz de respuestas o \(\delta y\) en el
vector, respectivamente,
\begin{align}
 \|\widehat H-H\|&\leq\frac{\|E\|}{\sqrt{589609}},
 \label{md:antmem:k:error-matriz}\\
 \|\widehat h-h\|&\leq\frac{\|\delta y\|}{\sqrt{589609}},\qquad
 \|\widehat H-H\|\leq\sqrt{\frac{12}{589609}}\,\|\delta y\|.
 \label{md:antmem:k:error-filtro}
\end{align}
\end{theorem}
\begin{proof}
La primera igualdad es \(Y=HO_K\). Si \(H\) conmuta con \(S\),
entonces \(HS^jK=S^jy\), de modo que \(Y=O_y\). La conmutación con
el ciclo determina todas las columnas de \(H\) desde una de ellas;
las doce potencias \(I,S,\ldots,S^{11}\) son independientes y generan
ese espacio de operadores. Por tanto \(y=O_Kh\), lo que da la
segunda fórmula. Las dos primeras cotas usan
\(\|O_K^{-1}\|=1/\sqrt{589609}\). La última usa
\(\|O_{\delta y}\|\leq\|O_{\delta y}\|_{\rm F}
=\sqrt{12}\|\delta y\|\).
\end{proof}

Una respuesta significa aquí un vector de doce componentes, no una
medición escalar. Como ejemplo, si \(H=2I+3S-S^5\), la respuesta
\(y=HK\) devuelve exactamente los coeficientes \(2,3,-1\) en sus
posiciones mediante \eqref{md:antmem:k:filtro}; el inversor necesita \(y,K,S\),
no los coeficientes del filtro. En cambio, el operador no nulo cuya primera
fila es \((543,-234,0,\ldots,0)\) y cuyas restantes filas son nulas
anula \(K\). Una sola respuesta no identifica un operador arbitrario.

La función de referencia admite una generalización exacta. En un ciclo de
\(n\) posiciones, \(H\mapsto Hg\) identifica todos los operadores
que conmutan con el ciclo si y sólo si
\(O_g=(g,Sg,\ldots,S^{n-1}g)\) es invertible. La suficiencia es la
prueba anterior. Si \(O_g\) es singular, un vector no nulo de su núcleo
da un polinomio de grado menor que \(n\) que anula \(g\); el operador
es no nulo porque las potencias del ciclo son independientes. El impulso
unitario también tiene órbita ortonormal y sirve de referencia. La
consecuencia específica para \(K\) es que el mismo objeto ya usado en
clausura e incidencia satisface esa condición, sin haber sido elegido de
nuevo para optimizar el ensayo. No se deduce una identificación de
dinámicas no lineales ni de estados ocultos arbitrarios; una aplicación a
una linealización necesita que ésta se haya definido y justificado.

\begin{corollary}[Recuperación de holonomía relativa desde el complemento]
\label{md:antmem:k:holonomia}
Sean \(H_B\) conocido e invertible y \(H_C\) otro transporte del
mismo tipo. El complemento
\[
 D_\gamma=\frac{\sqrt8}{9}(H_C-H_B)
\]
determina
\[
 H_B^{-1}H_C=I+\frac9{\sqrt8}H_B^{-1}D_\gamma.
\]
Si \(D_\gamma\) conmuta con \(S\), en particular si lo hacen ambos
transportes, basta \(y=D_\gamma K\) para recuperar
\begin{equation}
 H_B^{-1}H_C
 =I+\frac9{\sqrt8}H_B^{-1}O_yO_K^{-1}.
 \label{md:antmem:k:holonomia-una-respuesta}
\end{equation}
El error en esta reconstrucción está acotado por
\[
 \frac9{\sqrt8}\|H_B^{-1}\|
 \sqrt{\frac{12}{589609}}\,\|\delta y\|.
\]
\end{corollary}
\begin{proof}
Se despeja \(H_C\) en la definición del complemento y se aplica
\cref{md:antmem:k:operadores}. La cota resulta de la submultiplicatividad.
\end{proof}

Para \(H_B=S^3\) y \(H_C=S^4\), el dato racional
\(y/\sqrt8=(S^4-S^3)K/9\) permite reconstruir la holonomía relativa
\(S\) sin proporcionar \(H_C\) al inversor. Si \(H_B\) es
unitario, el factor de error es aproximadamente \(0.01435510\). Sin
conmutación se mantienen las doce respuestas generales, no una hipótesis
ficticia de simetría. La selección de un lazo físico y su representación
siguen perteneciendo a su dinámica generadora; este protocolo no las
sustituye.

\subsubsection{Consecuencia de la referencia marcada para dos orientaciones}
\label{md:antmem:orientaciones}
La consecuencia utilizada en las notas del cuaderno conserva la tabla
anterior como antecedente. Para un operador real \(H\) que conmuta con
\(S\), sean \(y_+=HK\) e \(y_-=H^{\mathsf T}K\).
Sus doce intensidades de Fourier coinciden, pero \(y_+=y_-\) si y sólo
si \(H=H^{\mathsf T}\). Cada respuesta vectorial completa, junto con
\(K\) y la orientación de \(S\), recupera su operador.
\begin{proof}
El operador \(H\) es circulante; sus multiplicadores de Fourier son
\(\lambda_m\), y los de \(H^{\mathsf T}\) son sus conjugados.
Por ello ambas respuestas tienen intensidades
\(|\lambda_m|^2|\widehat K_m|^2\).
Si sus vectores coinciden, \((H-H^{\mathsf T})K=0\).
La conmutación implica que \(H-H^{\mathsf T}\) anula también \(S^jK\)
para todo \(j\). Esos vectores forman una base, así que
\(H-H^{\mathsf T}=0\). La recuperación es
\(H=O_{y_+}O_K^{-1}\), y análogamente para \(H^{\mathsf T}\).
\end{proof}
En particular, los desplazamientos \(S\) y \(S^{-1}\) dan respuestas
de igual espectro de potencia, pero distintas. Las autocorrelaciones
exactas de la tabla anterior dan la identidad ya utilizada en el cuaderno:
\begin{equation}
 \|SK-S^{-1}K\|^2
 =2\bigl(\|K\|^2-\langle K,S^2K\rangle\bigr)
 =2\,175\,744>0.
 \label{md:antmem:distancia-orientada}
\end{equation}
Este número es una consecuencia de las autocorrelaciones, no un dato
de partida adicional. El control racional del cuaderno recupera de ambas
respuestas los coeficientes \(e_1\) y \(e_{11}\), respectivamente, por
inversión de \(O_K\). Una intensidad espectral o un balance escalar no
sustituye la respuesta vectorial.
Este corolario concierne al portador cíclico de doce componentes:
no identifica \(S\) con la dinámica enriquecida completa ni convierte
por sí solo \(S\leftrightarrow S^{-1}\) en conjugación física de carga.
Para \(H\) general, la transposición tampoco significa inversión temporal;
en el caso \(S\), sí coincide con \(S^{-1}\).

\subsection{Acoplamiento reversible y transporte no estacionario}
\label{md:antmem:acoplamiento}
\subsubsection{Extensión reversible del balance entre dos formas}

Sean \(B,C:V\to V'\) isomorfismos que conservan una forma no degenerada \(b\) entre dos realizaciones:

\[
B^*b'B=C^*b'C=b.
\]

La forma puede ser una métrica positiva o una forma alternante. En espacios reales, \(*\) denota transposición respecto de las coordenadas elegidas; la igualdad expresa el transporte de la forma entre las realizaciones especificadas. En espacios de Hilbert, los operadores y sus inversos son acotados.

Con \(s=\sqrt8\) y \(Q_V=\frac13\left(\begin{smallmatrix}sI&-I\\I&sI\end{smallmatrix}\right)\), definimos

\begin{equation}
\boxed{
\mathcal U(B,C)=Q_{V'}^*\operatorname{diag}(B,C)Q_V
=\frac19\begin{pmatrix}8B+C&\sqrt8(C-B)\\\sqrt8(C-B)&B+8C\end{pmatrix}.}
\label{md:antmem:acop:eq:1}
\end{equation}

Escribimos \(T=(8B+C)/9\), \(D=\sqrt8(C-B)/9\) y \(R=(B+8C)/9\). La actualización completa es

\begin{equation}
\binom{x'}{m'}=\begin{pmatrix}T&D\\D&R\end{pmatrix}\binom{x}{m}.
\label{md:antmem:acop:eq:2}
\end{equation}

\begin{theorem}
La aplicación \eqref{md:antmem:acop:eq:1} es reversible y conserva la forma \(b\oplus b\):

\begin{equation}
\mathcal U(B,C)^*(b'\oplus b')\mathcal U(B,C)=b\oplus b,
\quad\mathcal U(B,C)^{-1}=\mathcal U(B^{-1},C^{-1}).
\label{md:antmem:acop:eq:3}
\end{equation}
\end{theorem}

\begin{proof}
Se cumple \(Q^*Q=I\) por \(8+1=9\). Sus bloques escalares hacen que \(Q\) preserve \(b\oplus b\) para cualquier forma \(b\). La aplicación diagonal conserva las formas por hipótesis. La composición prueba \eqref{md:antmem:acop:eq:3}, y la inversión de los tres factores da la inversa declarada.
\end{proof}

Para memoria entrante \(m=0\), se recupera el balance \(b=T^*b'T+D^*b'D\). La segunda columna determina la respuesta a memoria entrante arbitraria. La conservación de normas positivas y la conservación de área orientada son las especializaciones respectivas de \eqref{md:antmem:acop:eq:3}.

\paragraph{Realización colectiva en las nueve copias}

En \(V^9\), sea \(W_E=\operatorname{diag}(I,\ldots,I,E)\), con \(E\) en la novena posición. Sean \(A(a,b)=(a/\sqrt8,\ldots,a/\sqrt8,b)\) y \(L=AQ\). La inclusión uniforme es \(Jx=(x,\ldots,x)/3\). Entonces \(L(x,0)=Jx\) y

\begin{equation}
W_E L=L\mathcal U(I,E).
\label{md:antmem:acop:eq:4}
\end{equation}

\begin{proof}
\(W_EA=A\operatorname{diag}(I,E)\); multiplicar por \(Q\) y usar \(AQ=L\) da \eqref{md:antmem:acop:eq:4}. La ortogonalidad de la media de las ocho primeras posiciones con sus siete diferencias muestra además que el complemento queda fijo bajo \(W_E\).
\end{proof}

Sea \(P_{\rm cyc}(v_0,\ldots,v_8)=(v_8,v_0,\ldots,v_7)\). El paso cíclico original es \(S_E=P_{\rm cyc}W_E\). Por ello \(S_EL=(P_{\rm cyc}L)U(I,E)\), con coordenatizaciones de entrada y salida distintas. La reducción colectiva conserva el patrón correspondiente: \(S_E^9=I_9\otimes E\), mientras \(W_E^9=\operatorname{diag}(I_8,E^9)\). Cuando un lector utiliza las nueve etiquetas individuales, éstas se conservan en su representación original.

\subsubsection{Composición exacta y reentrada de memoria}

Para \(B_1,C_1:V_0\to V_1\) y \(B_2,C_2:V_1\to V_2\),

\begin{equation}
\boxed{\mathcal U(B_2,C_2)\mathcal U(B_1,C_1)
=\mathcal U(B_2B_1,C_2C_1).}
\label{md:antmem:acop:eq:5}
\end{equation}

\begin{proof}
En \eqref{md:antmem:acop:eq:1} se cancelan \(Q_{V_1}Q_{V_1}^*\), y los dos bloques diagonales se multiplican en el orden indicado. La inducción extiende la fórmula a cualquier número finito de transportes, conservando su orden.
\end{proof}

El bloque de lectura contiene la identidad

\begin{equation}
\boxed{T_2T_1+D_2D_1=\frac{8B_2B_1+C_2C_1}{9}.}
\label{md:antmem:acop:eq:6}
\end{equation}

Desde \((x,0)\), la primera etapa produce \((T_1x,D_1x)\). La segunda lectura es \(T_2T_1x+D_2D_1x\). El segundo término es exactamente la contribución de la memoria reintroducida. Al almacenar \(D_1x\) fuera del siguiente acoplamiento y continuar sólo con \(T_1x\) se obtiene el procedimiento de registros separados, cuya lectura es \(T_2T_1x\). La igualdad \eqref{md:antmem:acop:eq:6} precisa la diferencia entre ambos procedimientos.

La realización \(U\) conserva el estado del acoplamiento y los productos ordenados \(B_N\cdots B_1\), \(C_N\cdots C_1\). La conservación de todos los prefijos incluye además el registro de emisiones de HMT, que mantiene cada factor junto con su orden.

\paragraph{Ecuación covariante de incremento}

Para el registro separado \(x_{n+1}=T_nx_n\), \(m_n=D_nx_n\), se tiene

\begin{equation}
x_{n+1}-B_nx_n=m_n/\sqrt8.
\label{md:antmem:acop:eq:7}
\end{equation}

Con \(G_0=I\) y \(G_{n+1}=B_nG_n\), la telescopía da

\begin{equation}
x_0=G_N^{-1}x_N-\frac1{\sqrt8}\sum_{n<N}G_{n+1}^{-1}m_n.
\label{md:antmem:acop:eq:8}
\end{equation}

La memoria es el incremento respecto del transporte entre formas \(B_n\), con la normalización declarada. La igualdad \eqref{md:antmem:acop:eq:8} es una reconstrucción finita algebraica; su paso a una serie infinita requiere convergencia. El teorema siguiente construye la inversa estable mediante el límite de operadores positivos.

\subsubsection{Memoria de una sucesión no estacionaria}

En una realización de Hilbert \(H\), sean \(C_n\) unitarios, con su orden de composición declarado. Escribimos

\[
T_n=(8I+C_n)/9,\quad D_n=\sqrt8(C_n-I)/9,
\quad P_0=I,\quad P_{n+1}=T_nP_n,\quad A_N=P_N^*P_N.
\]

El índice \(N\) cuenta compresiones con almacenamiento separado, a diferencia del transporte completo con reentrada. Para métricas variables con \(B_n\) isomorfismos isométricos, la conjugación por \(G_n\) lleva los transportes a un Hilbert fijo y produce este mismo caso, con \(\widetilde C_n=G_{n+1}^{-1}C_nG_n\). La afirmación utiliza las métricas de las fibras y las isometrías declaradas entre ellas.

\begin{theorem}[Terminal positivo y recuperación no estacionaria]
\label{md:antmem:terminal-no-estacionario}
Existe un operador positivo \(0\leq A_\infty\leq I\), límite fuerte de \(A_N\), y

\begin{equation}
\boxed{I=A_\infty+\sum_{n\ge0}P_n^*D_n^*D_nP_n.}
\label{md:antmem:acop:eq:9}
\end{equation}

La aplicación

\begin{equation}
\boxed{\mathcal Vv=(A_\infty^{1/2}v,(D_nP_nv)_{n\ge0})}
\label{md:antmem:acop:eq:10}
\end{equation}

es isométrica y tiene inversa izquierda estable

\begin{equation}
\mathcal V^*(a,(m_n))=A_\infty^{1/2}a+
\sum_{n\ge0}P_n^*D_n^*m_n,\qquad\|\mathcal V^*\|=1.
\label{md:antmem:acop:eq:11}
\end{equation}

\end{theorem}
\begin{proof}
La identidad local \(T_n^*T_n+D_n^*D_n=I\) da
\(A_N-A_{N+1}=P_N^*D_N^*D_NP_N\geq0\). Las formas cuadráticas de
la sucesión positiva decreciente convergen a la de un operador positivo
\(A_\infty\). Para \(B_N=A_N-A_\infty\), la desigualdad
\(0\leq B_N^2\leq B_N\) convierte esa convergencia en fuerte.
La telescopía finita y el límite prueban \eqref{md:antmem:acop:eq:9}. Evaluar sobre \(v\)
prueba \eqref{md:antmem:acop:eq:10}. Su adjunto es \eqref{md:antmem:acop:eq:11}; la serie converge en norma porque
las truncaciones de un registro cuadrado-sumable convergen y el adjunto es acotado.
\end{proof}

La reconstrucción con el terminal límite y \(N\) registros tiene error
exacto \((A_N-A_\infty)v\). El vector \(A_\infty^{1/2}v\) es la
realización positiva canónica del terminal de norma. Su existencia procede
del límite de Gram, con independencia de la convergencia vectorial de \(P_Nv\).

\subsubsection{Tres propiedades diferentes: agotamiento, distinción y estabilidad}

Sea \(Mv=(D_nP_nv)_n\). De \eqref{md:antmem:acop:eq:9}, \(M^*M=I-A_\infty\). Por tanto:

\begin{enumerate}
\item Toda la norma de \(v\) pasa al registro exactamente cuando \(A_\infty v=0\).
\item El registro distingue cada estado exactamente cuando \(\ker(I-A_\infty)=\{0\}\).
\item La recuperación desde memoria sola tiene inversa acotada exactamente cuando
\(I-A_\infty\geq\varepsilon I\) para algún \(\varepsilon>0\).
En ese caso una inversa izquierda es \((I-A_\infty)^{-1}M^*\).
\end{enumerate}

Además,

\begin{equation}
\boxed{\ker M=\bigcap_{n\ge0}\operatorname{Fix}(C_n).}
\label{md:antmem:acop:eq:12}
\end{equation}

\begin{proof}
Si \(Mv=0\), entonces \(D_0v=0\), de donde \(C_0v=v\)
y \(P_1v=v\). Inductivamente, \(D_nP_nv=0\) da \(C_nv=v\) y
\(P_{n+1}v=v\). El recíproco es inmediato. Las tres equivalencias se
deducen del Gram \(M^*M\) y de la caracterización de los operadores
acotados inferiormente.
\end{proof}

\paragraph{Ejemplo exacto que separa las tres propiedades}

En el ejemplo algebraico \(C_0=-I\) y \(C_n=I\) para \(n\geq1\),

\begin{equation}
A_\infty=\frac{49}{81}I,\qquad M^*M=\frac{32}{81}I,
\qquad m_0=-\frac{2\sqrt8}{9}v,
\qquad v=-\frac9{2\sqrt8}m_0.
\label{md:antmem:acop:eq:13}
\end{equation}

La memoria contiene \(32/81\) de la norma al cuadrado y permite recuperar el vector completo. El terminal retiene \(49/81\) y es un operador positivo distinto de un proyector. La recuperabilidad depende del núcleo y del condicionamiento del lector. Lectura y memoria constituyen componentes correlacionadas de una isometría.

Este ejemplo separa las tres consecuencias lógicas del teorema. La
aplicación a una trayectoria física conserva además la selección y la
procedencia efectiva de sus \(C_n\) mediante el TPK.

La letra \(\eta\) de la realización gaussiana siguiente designa la
deformación de la elipse, es decir, el parámetro \(\zeta\) de los
desarrollos de transporte; no es el coeficiente de acción
\(\eta_{\rm ret}\). Se mantiene una sección positiva \(\hbar_a\) y
una misma unidad para las dos coordenadas del cociente angular.
Los semiejes conservados son
\[
 a_S=\sqrt{2\hbar_a}\,e^{\eta/2},\qquad
 b_S=\sqrt{2\hbar_a}\,e^{-\eta/2},\qquad
 \eta=\operatorname{artanh}(C^*/A).
\]

\subsection{Covarianza canónica y reciprocidad de la elipse de acción}
\label{md:antmem:covarianza-accion}

La construcción precedente determina dos invariantes complementarios:
el producto de los semiejes fija el área de acción y su cociente conserva
la anisotropía del par angular. Precisamos ahora la realización canónica
en la que esos semiejes son dos desviaciones estándar. La estructura de
partida es la elipse ya obtenida de las representaciones angulares y de
acción de APP--TRIT--TPK; la realización operatoria se establece después
de esa construcción.

En este apartado escribimos \(\eta=\eta_{\rm el}
=\operatorname{artanh}(C^*/A)\), con \(|C^*/A|<1\) y
\(\hbar_a>0\). El cociente angular utiliza una misma unidad para sus dos
componentes. En las coordenadas equilibradas \((Q,P)\), ambas de
dimensión raíz de acción, definimos
\begin{equation}
 V_\eta:=\frac14
 \begin{pmatrix}a_S^2&0\\0&b_S^2\end{pmatrix}
 =\frac{\hbar_a}{2}
 \begin{pmatrix}e^\eta&0\\0&e^{-\eta}\end{pmatrix}.
 \label{md:antmem:covarianza-geometrica}
\end{equation}
Esta matriz positiva tiene entradas de dimensión acción. Su interpretación
como covarianza cuántica queda especificada por la realización siguiente.

\begin{proposition}[Realización gaussiana de la covarianza de acción]
\label{md:antmem:covarianza-gaussiana}
Considérese la representación canónica en \(L^2(\mathbb R,dQ)\), con
\(\widehat Qf(Q)=Qf(Q)\) y
\(\widehat Pf(Q)=-i\hbar_a f'(Q)\), inicialmente sobre el espacio de
Schwartz \(\mathcal S(\mathbb R)\). El estado normalizado
\begin{equation}
 \psi_\eta(Q)=\bigl(\pi\hbar_a e^\eta\bigr)^{-1/4}
 \exp\!\left(-\frac{Q^2}{2\hbar_a e^\eta}\right)
 \label{md:antmem:gaussiana-accion}
\end{equation}
tiene medias nulas y matriz de covarianza \(V_\eta\). En particular,
\begin{equation}
 (\Delta Q)^2=\frac{\hbar_a e^\eta}{2},\qquad
 (\Delta P)^2=\frac{\hbar_a e^{-\eta}}{2},\qquad
 \operatorname{Cov}(Q,P)=0,\qquad
 \det V_\eta=\frac{\hbar_a^2}{4}.
 \label{md:antmem:varianzas-accion}
\end{equation}
Esta realización satura la desigualdad de Robertson--Schrödinger.
\end{proposition}

\begin{proof}
Los operadores anteriores son esencialmente autoadjuntos sobre
\(\mathcal S(\mathbb R)\), que constituye un núcleo común invariante;
sobre este espacio se cumple
\([\widehat Q,\widehat P]=i\hbar_a I\). Para
\(\widehat Y=(\widehat Q,\widehat P)\), la covarianza simétrica es
\[
 V_{jk}=\frac12\left\langle
 (\widehat Y_j-\langle\widehat Y_j\rangle)
 (\widehat Y_k-\langle\widehat Y_k\rangle)
 +(\widehat Y_k-\langle\widehat Y_k\rangle)
 (\widehat Y_j-\langle\widehat Y_j\rangle)
 \right\rangle.
\]
Con \(s=\hbar_a e^\eta>0\), sea
\(I_s=\int_{\mathbb R}e^{-Q^2/s}\,dQ\). El producto de dos
integrales iguales, evaluado en coordenadas polares, da
\(I_s^2=2\pi\int_0^\infty r e^{-r^2/s}\,dr=\pi s\).
Además, la integral de
\(\frac{d}{dQ}(Qe^{-Q^2/s})\) es cero, por el decaimiento en ambos
extremos; de ahí
\(\int_{\mathbb R}Q^2e^{-Q^2/s}\,dQ=sI_s/2\).
Estas identidades y la paridad prueban la normalización, la media nula y
\(\langle\widehat Q^2\rangle=s/2\). Además,
\(\widehat P\psi_\eta=i e^{-\eta}Q\psi_\eta\).
Por tanto, \(\langle\widehat P\rangle=0\),
\(\|\widehat P\psi_\eta\|^2=\hbar_a e^{-\eta}/2\) y
\(\operatorname{Re}\langle\widehat Q\psi_\eta,
\widehat P\psi_\eta\rangle=0\).

Para cualquier estado normalizado de ese dominio, Cauchy--Schwarz aplicada
a \((\widehat Q-\langle\widehat Q\rangle)\psi\) y
\((\widehat P-\langle\widehat P\rangle)\psi\) da
\[
 (\Delta Q)^2(\Delta P)^2
 \geq \operatorname{Cov}(Q,P)^2
       +\frac14\left|\langle[\widehat Q,\widehat P]\rangle\right|^2
 =\operatorname{Cov}(Q,P)^2+\frac{\hbar_a^2}{4}.
\]
Las varianzas calculadas alcanzan la igualdad. Equivalentemente, el
operador
\[
 a_\eta=\frac{e^{-\eta/2}\widehat Q
                  +i e^{\eta/2}\widehat P}{\sqrt{2\hbar_a}}
\]
satisface \([a_\eta,a_\eta^\dagger]=I\) y
\(a_\eta\psi_\eta=0\), por sustitución directa.
\end{proof}

\subsection{Covarianza, forma normal y espectro gaussiano}
\label{md:antmem:gauss:gaussianas-espectro}

La sección positiva de acción \(\hbar_a\) y la realización canónica
anterior determinan la normalización empleada aquí. Se demuestra la
correspondencia entre covarianza, espectro, pureza y entropía para un
número finito de modos; el espacio de Fock conserva todas sus ocupaciones.
La reducción simpléctica y las fórmulas gaussianas son resultados
establecidos que se reproducen para hacer autónoma su aplicación al
transporte de memoria.

\subsubsection{Coordenadas canónicas y función característica}
En \(L^2(\mathbb R^d)\), sean \(\widehat q_j=\widehat Q_j/\sqrt{\hbar_a}\),
\(\widehat p_j=\widehat P_j/\sqrt{\hbar_a}\) y
\(\widehat z=(\widehat q_1,\widehat p_1,\ldots,\widehat q_d,\widehat p_d)\).
Sobre Schwartz,
\[
 [\widehat z_j,\widehat z_k]=iJ_{jk}I,\qquad
 J=\bigoplus_{j=1}^d\begin{pmatrix}0&1\\-1&0\end{pmatrix}.
\]
Esta convención fija el signo de la forma canónica. Los operadores de
Weyl \(\mathcal W(\xi)=\exp(i\xi^{\mathsf T}\widehat z)\) se definen
por las clausuras autoadjuntas de las combinaciones lineales indicadas.
Para un estado centrado \(\rho\) de segundos momentos finitos, la
covarianza normalizada y su función característica son
\[
 W_{jk}=\operatorname{Tr}\rho(\widehat z_j\widehat z_k+\widehat z_k\widehat z_j),
 \qquad \chi_\rho(\xi)=\operatorname{Tr}(\rho\mathcal W(\xi)),\qquad
 V=\frac{\hbar_a}{2}W.
\]
El estado es gaussiano centrado cuando
\begin{equation}
 \chi_\rho(\xi)=\exp(-\xi^{\mathsf T}W\xi/4).
 \label{md:antmem:gauss:eq:caracteristica-gaussiana}
\end{equation}
La esperanza de \((c^{\mathsf T}\widehat z)^\dagger
(c^{\mathsf T}\widehat z)\) es \(c^\dagger(W+iJ)c/2\); así
\(W+iJ\geq0\). Además \(W>0\): si \(v\) es real y
\(v^{\mathsf T}Wv=0\), la positividad implica \((W+iJ)v=0\),
por tanto \(Jv=0\) y \(v=0\).

\begin{lemma}[Unicidad de la función característica]
\label{md:antmem:gauss:lem:unicidad-weyl}
Dos operadores de clase traza con la misma función característica son iguales.
\end{lemma}
\begin{proof}
Reagrupando posiciones y momentos como \(\xi=(u,v)\),
\[
 (\mathcal W(u,v)f)(x)=e^{iu\cdot(x+v/2)}f(x+v).
\]
Para un operador de rango finito con vectores de Schwartz y núcleo \(K_A\),
\[
 \operatorname{Tr}(A\mathcal W(u,v))=\int K_A(x+v/2,x-v/2)e^{iu\cdot x}\,dx.
\]
Plancherel y el cambio de variables de jacobiano absoluto uno dan
\[
 \|A\|_{\rm HS}^2=(2\pi)^{-d}\int_{\mathbb R^{2d}}
 |\operatorname{Tr}(A\mathcal W(\xi))|^2\,d\xi.
\]
Los operadores anteriores son densos en clase traza. La convergencia
en esa norma implica convergencia en Hilbert--Schmidt y convergencia
uniforme de las funciones características, porque su diferencia está
acotada por \(\|A-B\|_1\). La igualdad se extiende en \(L^2\).
Aplicarla a la diferencia de los operadores prueba la afirmación.
\end{proof}

\subsubsection{Reducción simpléctica de la forma positiva}
\begin{theorem}[Forma normal de Williamson]
\label{md:antmem:gauss:thm:williamson}
Para una matriz real simétrica \(W>0\) de orden \(2d\), existen
\(S\in\operatorname{Sp}(2d,\mathbb R)\) y \(\nu_j>0\) tales que
\[
 W=SDS^{\mathsf T},\qquad D=\bigoplus_j\nu_jI_2.
\]
Los autovalores de \(JW\) son \(\pm i\nu_j\). La condición
\(W+iJ\geq0\) equivale a \(\nu_j\geq1\) para todo \(j\).
\end{theorem}
\begin{proof}
La matriz \(B=W^{1/2}JW^{1/2}\) es antisimétrica e invertible.
Si \(-B^2u=\nu^2u\) y \(\|u\|=1\), el vector
\(v=-Bu/\nu\) es unitario y ortogonal a \(u\); cumple
\(Bu=-\nu v\), \(Bv=\nu u\). El complemento de esta pareja es
invariante. La iteración produce una matriz ortogonal \(O\) con
\(O^{\mathsf T}BO=JD\). Definiendo \(S=W^{1/2}OD^{-1/2}\),
\[
 S^{\mathsf T}JS=J,\qquad SDS^{\mathsf T}=W.
\]
La semejanza \(W^{1/2}(JW)W^{-1/2}=B\) determina los valores
\(\nu_j\), con multiplicidad. La congruencia por \(S^{-1}\)
transforma \(W+iJ\) en \(D+iJ\); cada bloque tiene autovalores
\(\nu_j+1\) y \(\nu_j-1\).
\end{proof}

\begin{lemma}[Implementación unitaria en un número finito de modos]
\label{md:antmem:gauss:lem:implementacion-simplectica}
Toda matriz simpléctica \(S\) admite un unitario \(U_S\) que satisface
\(U_S^\dagger\mathcal W(\xi)U_S=\mathcal W(S^{\mathsf T}\xi)\).
La conjugación \(\rho\mapsto U_S\rho U_S^\dagger\) transporta
la covarianza por \(W\mapsto SWS^{\mathsf T}\).
\end{lemma}
\begin{proof}
En la descomposición polar \(S=OP\), la identidad
\(J^{-1}(S^{\mathsf T}S)J=(S^{\mathsf T}S)^{-1}\), aplicada por
cálculo espectral a la raíz positiva, prueba \(PJP=J\).
Entonces \(O\) es ortogonal simpléctica. Los autoespacios de \(P\)
se emparejan mediante \(J\) con autovalores \(\lambda,\lambda^{-1}\);
el de autovalor uno es \(J\)-invariante. Por ello una matriz
ortogonal simpléctica \(K\) reduce \(P\) a
\(\bigoplus_j\operatorname{diag}(e^{r_j},e^{-r_j})\).

La dilatación se implementa por
\[
 (\mathcal D_rf)(x)=e^{-\frac12\sum_jr_j}
 f(e^{-r_1}x_1,\ldots,e^{-r_d}x_d).
\]
Un cambio de variables prueba su unitariedad y la conjugación de
\((\widehat q_j,\widehat p_j)\) a
\((e^{r_j}\widehat q_j,e^{-r_j}\widehat p_j)\).
Una transformación ortogonal simpléctica conmuta con \(J\) y
corresponde a una matriz \(u\in U(d)\) sobre
\(a_j=(\widehat q_j+i\widehat p_j)/\sqrt2\). Su segunda
cuantización \(\Gamma(u)=\bigoplus_{n\geq0}u^{\otimes_s n}\)
es unitaria en \(\mathcal F_s(\mathbb C^d)\), y la identidad sobre
productos simétricos finitos da la transformación de los \(a_j\).

La base de Hermite identifica este espacio con \(L^2(\mathbb R^d)\).
Para justificar su completitud, si \(f\) es ortogonal a todos los
polinomios multiplicados por \(e^{-|x|^2/2}\), la función entera
\(F(z)=\int f(x)e^{-|x|^2/2}e^{z\cdot x}\,dx\) tiene nulas
todas sus derivadas en cero. Así \(F=0\); restringiendo a
\(z=i\xi\), la unicidad de Fourier da \(f=0\). Las
normalizaciones y la ortogonalidad resultan de creación y aniquilación.
La composición de las implementaciones anteriores produce \(U_S\).
La igualdad de Weyl implica
\(\chi_{U_S\rho U_S^\dagger}(\xi)=\chi_\rho(S^{\mathsf T}\xi)\),
que demuestra la ley de covarianza.
\end{proof}

\subsubsection{Cálculo del espectro de la densidad isotrópica}
En un modo, los estados coherentes normalizados son
\[
 |z\rangle=e^{-|z|^2/2}\sum_{n\geq0}\frac{z^n}{\sqrt{n!}}|n\rangle.
\]
Su función de posición,
\(\pi^{-1/4}\exp[-(x-q_0)^2/2+ip_0(x-q_0/2)]\), donde
\(q_0=\sqrt2\Re z\), \(p_0=\sqrt2\Im z\), da por integración
\[
 \langle z|\mathcal W(u,v)|z\rangle
 =e^{-(u^2+v^2)/4}e^{i(uq_0+vp_0)}.
\]
Para \(s>0\), la integral
\[
 \rho^{(s)}=\frac1{\pi s}\int_{\mathbb C} e^{-|z|^2/s}
 |z\rangle\langle z|\,d^2z
\]
converge en clase traza, es positiva y tiene traza uno. La integración
angular anula sus elementos fuera de la diagonal de ocupación; la
integración radial da
\[
 \langle n|\rho^{(s)}|n\rangle=\frac{s^n}{(1+s)^{n+1}},\qquad
 \chi_{\rho^{(s)}}(u,v)=e^{-(2s+1)(u^2+v^2)/4}.
\]
Con \(\nu=2s+1\), el único estado gaussiano de covarianza
\(\nu I_2\) es, por el lema~\ref{md:antmem:gauss:lem:unicidad-weyl},
\begin{equation}
 \rho_\nu=(1-t)\sum_{n\geq0}t^n|n\rangle\langle n|,
 \qquad t=\frac{\nu-1}{\nu+1}.
 \label{md:antmem:gauss:eq:estado-geometrico}
\end{equation}
El caso \(\nu=1\) es el proyector sobre el vacío.

\begin{theorem}[Espectro, pureza y entropía]
\label{md:antmem:gauss:thm:espectro-gaussiano}
Un estado gaussiano centrado con covarianza \(W\) y valores
simplécticos \(\nu_1,\ldots,\nu_d\) es unitariamente equivalente
a \(\bigotimes_j\rho_{\nu_j}\). Sus autovalores son
\[
 \lambda_{\mathbf n}=\prod_j(1-t_j)t_j^{n_j},\qquad
 t_j=\frac{\nu_j-1}{\nu_j+1},\quad \mathbf n\in\mathbb N^d.
\]
Con \(0\log0=0\),
\begin{align}
 \operatorname{Tr}\rho^2&=\prod_j\nu_j^{-1}=(\det W)^{-1/2},
 \label{md:antmem:gauss:eq:pureza-general}\\
 -\operatorname{Tr}(\rho\log\rho)&=\sum_j
 \left[\frac{\nu_j+1}{2}\log\frac{\nu_j+1}{2}
 -\frac{\nu_j-1}{2}\log\frac{\nu_j-1}{2}\right].
 \label{md:antmem:gauss:eq:entropia-general}
\end{align}
\end{theorem}
\begin{proof}
Williamson da \(W=SDS^{\mathsf T}\). El producto de las densidades
isotrópicas tiene función característica \(e^{-\xi^{\mathsf T}D\xi/4}\).
La implementación unitaria y la unicidad de Weyl prueban la
equivalencia y la fórmula espectral. En un modo,
\[
 \sum_{n\geq0}(1-t)^2t^{2n}=\frac{1-t}{1+t}=\nu^{-1},
\]
y
\[
 -\sum_{n\geq0}(1-t)t^n\log((1-t)t^n)
 =-\log(1-t)-\frac{t}{1-t}\log t.
\]
La sustitución \(t=(\nu-1)/(\nu+1)\) prueba las fórmulas; el
límite en \(t=0\) es cero. En el producto finito, Tonelli justifica
la suma de los términos de entropía no negativos, y
\(\det W=\prod_j\nu_j^2\) da la expresión determinantal.
\end{proof}

Si el estado es la reducción de un vector puro bipartito, sus
autovalores son los cuadrados de los coeficientes de Schmidt.
La descomposición singular del operador de Hilbert--Schmidt definido
por ese vector demuestra la igualdad de espectros no nulos de ambas
reducciones. La entropía calculada es entonces la de entrelazamiento.

\subsubsection{Preparación y transporte de las dos covarianzas}
La aplicación posterior utiliza dos entradas gaussianas puras idénticas.
En general, para \(M\in\operatorname{Sp}(2d,\mathbb R)\), su covarianza es
\(V_0=(\hbar_a/2)MM^{\mathsf T}\), y la entrada conjunta tiene
covarianza \(V_0\oplus V_0\). El transporte físico se obtiene
conjugando \(U_H\) por \(M\oplus M\). En la coordenatización equilibrada,
la covarianza normalizada de entrada es \(I_{4d}\); con
\[
 T=\frac{8I+H}{9},\qquad D=\frac{\sqrt8(H-I)}9,
\]
la primera reducción tiene
\[
 W=TT^{\mathsf T}+DD^{\mathsf T}=\frac{8I+HH^{\mathsf T}}9.
\]
Su función característica se obtiene anulando los argumentos del
segundo grupo de modos. Por ello permanece gaussiana y los teoremas
anteriores calculan todo su espectro, sin truncar el registro de ocupación.

\subsection{Prolongación del transporte de memoria: realización cuántica}

La realización cuántica utiliza el acoplamiento reversible de la sección~\ref{md:antmem:acoplamiento}. Consideramos primero un plano canónico real con matriz alternante \(\Omega=\left(\begin{smallmatrix}0&1\\-1&0\end{smallmatrix}\right)\). La partición nonádica determina

\begin{equation}
Q=\frac13\begin{pmatrix}\sqrt8 I&-I\\I&\sqrt8 I\end{pmatrix},\qquad
U_H=Q^{\mathsf T}\operatorname{diag}(I,H)Q
=\frac19\begin{pmatrix}8I+H&\sqrt8(H-I)\\
\sqrt8(H-I)&I+8H\end{pmatrix}.
\label{md:antmem:cuantica:eq:5.1}
\end{equation}

El operador \(H\in\operatorname{Sp}(2,\mathbb R)\) transporta ese plano. La matriz \(Q\) conserva el producto euclídeo y la forma \(\Omega\oplus\Omega\); por ello \(U_H\) es simpléctica. La primera componente define el sistema observado y la segunda el registro de memoria; la transformación completa actúa sobre ambos.

El lazo ordenado cuyas identidades se reproducen en
\ref{md:antmem:lazo-ep} es

\begin{equation}
C=\begin{pmatrix}0&-1\\1&-1\end{pmatrix},\quad
P=\begin{pmatrix}1&1\\0&1\end{pmatrix},\quad
H_n=C^{-1}P^{-n}CP^n
=\begin{pmatrix}1+n&n^2\\n&n^2-n+1\end{pmatrix}.
\label{md:antmem:cuantica:eq:5.2}
\end{equation}

Para cada \(n\in\mathbb Z\) se cumple \(\det H_n=1\). En particular, para \(n=1\),

\begin{equation}
H_1=\begin{pmatrix}2&1\\1&1\end{pmatrix}
=F_1^2,\quad F_1=\begin{pmatrix}1&1\\1&0\end{pmatrix}
=P F_{\rm av}P^{-1}.
\label{md:antmem:cuantica:eq:5.3}
\end{equation}

Su espectro es \(\{\varphi^2,\varphi^{-2}\}\), por el reconocimiento posterior del modo áureo ya generado. El entero \(n\) especifica el número orientado de iteraciones en el lazo elegido. La colección de transportes está determinada algebraicamente; su realización como protocolo físico conserva además los datos de preparación y de evolución correspondientes.

\subsubsection{Identidades algebraicas del lazo y reconocimiento posterior}
\label{md:antmem:lazo-ep}
Para conservar la prueba de las identidades que se utilizan, se reproduce
su composición matricial. En la misma coordenatización,
\[
 N_0=\begin{pmatrix}0&1\\0&0\end{pmatrix}.
\]
Sea
\[
 P=I+N_0=\begin{pmatrix}1&1\\0&1\end{pmatrix}.
\]
Como \(N_0^2=0\), se cumple \(P^n=I+nN_0\) para todo
\(n\in\mathbb Z\).
El lazo ordenado de cuatro transportes da

\begin{equation}
H_n=C^{-1}P^{-n}CP^n
=\begin{pmatrix}1+n&n^2\\n&n^2-n+1\end{pmatrix}.
\label{md:antmem:matriz:eq:28}
\end{equation}

Para \(n\ne0\), su diferencia normalizada es

\[
F_n=\frac{H_n-I}{n}
=\begin{pmatrix}1&n\\1&n-1\end{pmatrix}.
\]

Multiplicando esas matrices se obtiene, sin introducir valores espectrales,

\begin{equation}
\boxed{F_n^2=H_n,\qquad F_n^2-nF_n-I=0.}
\label{md:antmem:matriz:eq:29}
\end{equation}

Con \(H_B=I\), la memoria es
\(D_{\gamma,n}=(\sqrt8/9)nF_n\).
En consecuencia,

\begin{equation}
\boxed{H_n=\left(\frac9{n\sqrt8}D_{\gamma,n}\right)^2.}
\label{md:antmem:matriz:eq:30}
\end{equation}

La respuesta matricial de memoria reconstruye la autoescala
hiperbólica. Para \(n>0\), el reconocimiento espectral posterior
identifica

\begin{equation}
\rho_{{\rm ep},n}=\frac{n+\sqrt{n^2+4}}2,\quad
\operatorname{Spec}F_n=\{\rho_{{\rm ep},n},-\rho_{{\rm ep},n}^{-1}\},\quad
\operatorname{Spec}H_n=\{\rho_{{\rm ep},n}^2,\rho_{{\rm ep},n}^{-2}\}.
\label{md:antmem:matriz:eq:31}
\end{equation}

El caso \(n=1\) produce el polinomio áureo.
La relación con el operador
\(F_{\rm av}=\begin{pmatrix}0&1\\1&1\end{pmatrix}\)
generado por el calendario conserva el cambio
de coordenatización:

\begin{equation}
F_1=PF_{\rm av}P^{-1},\qquad
H_1=PF_{\rm av}^{\,2}P^{-1}.
\label{md:antmem:matriz:eq:32}
\end{equation}

El transporte \(P\), de determinante uno, entrelaza ambas
matrices y conserva \(\Omega\).

\subsubsection{Estado inicial, transporte y lectura}

La proposición~\ref{md:antmem:covarianza-gaussiana} construye la realización gaussiana de la elipse de acción. En la sección positiva \(\hbar_a>0\), con \(A>0\) y \(|C^*/A|<1\), su covarianza es

\begin{equation}
V_\eta=\frac{\hbar_a}{2}M_\eta M_\eta^{\mathsf T},\qquad
M_\eta=\operatorname{diag}(e^{\eta/2},e^{-\eta/2}),\qquad
\eta=\operatorname{artanh}(C^*/A).
\label{md:antmem:cuantica:eq:5.4}
\end{equation}

En la representación de Schrödinger sobre \(L^2(\mathbb R,dQ)\), con \(\widehat Q\) multiplicación por \(Q\) y \(\widehat P=-i\hbar_a\,d/dQ\), la relación \([\widehat Q,\widehat P]=i\hbar_a I\) vale en el espacio de Schwartz. El estado puro normalizado es

\[
\psi_\eta(Q)=(\pi\hbar_a e^\eta)^{-1/4}
\exp[-Q^2/(2\hbar_a e^\eta)].
\]

La preparación es el producto \(\psi_\eta\otimes\psi_\eta\). El transporte de \eqref{md:antmem:cuantica:eq:5.1} se expresa en la misma coordenatización de covarianza por

\begin{equation}
\widetilde U=(M_\eta\oplus M_\eta)U_H
(M_\eta^{-1}\oplus M_\eta^{-1}).
\label{md:antmem:cuantica:eq:5.5}
\end{equation}

El lema~\ref{md:antmem:gauss:lem:implementacion-simplectica} proporciona una implementación unitaria de esta matriz simpléctica en dos modos. Las implementaciones que difieren por una fase global producen la misma conjugación del estado. La covarianza se transforma por \(V\mapsto\widetilde U V\widetilde U^{\mathsf T}\); el estado conjunto permanece puro y su traza parcial sobre el segundo modo determina la lectura visible, conforme a la construcción de la sección~\ref{md:antmem:gauss:gaussianas-espectro}.

\begin{theorem}[Mezcla reducida por deformación relativa]

La covarianza de la lectura visible es

\begin{equation}
\boxed{V_{\rm vis}=
\frac{\hbar_a}{2}M_\eta
\frac{8I+HH^{\mathsf T}}9
M_\eta^{\mathsf T}.}
\label{md:antmem:cuantica:eq:5.6}
\end{equation}

Su valor simpléctico normalizado es

\begin{equation}
\boxed{\nu(H):=\frac{2\sqrt{\det V_{\rm vis}}}{\hbar_a}
=\sqrt{1+\frac8{81}\left[\operatorname{tr}(HH^{\mathsf T})-2\right]}.}
\label{md:antmem:cuantica:eq:5.7}
\end{equation}
\end{theorem}

\begin{proof}
Sean \(T=(8I+H)/9\) y \(D=\sqrt8(H-I)/9\). Los dos modos iniciales son independientes y tienen la misma covarianza. En la coordenatización equilibrada del estado inicial, la covarianza visible normalizada es \(TT^{\mathsf T}+DD^{\mathsf T}\). La expansión da

\[
\begin{aligned}
TT^{\mathsf T}+DD^{\mathsf T}
&=\frac{64I+8(H+H^{\mathsf T})+HH^{\mathsf T}
+8[HH^{\mathsf T}-H-H^{\mathsf T}+I]}{81}\\
&=\frac{8I+HH^{\mathsf T}}9.
\end{aligned}
\]

Esto demuestra \eqref{md:antmem:cuantica:eq:5.6}. Puesto que \(\det M_\eta=1\) y \(\det(HH^{\mathsf T})=1\),
\(\det(8I+HH^{\mathsf T})=65+8\operatorname{tr}(HH^{\mathsf T})\).
La división por \(81\) demuestra \eqref{md:antmem:cuantica:eq:5.7}. Los dos autovalores positivos de \(HH^{\mathsf T}\) son recíprocos; su suma es al menos dos. En consecuencia, \(\nu\geq1\), con igualdad exactamente cuando \(HH^{\mathsf T}=I\).
\end{proof}

La escala \(\hbar_a\) y la deformación \(\eta\) se conservan hasta efectuar el determinante. Su cancelación es covariante porque se transportan conjuntamente estado y operador. La expresión \(HH^{\mathsf T}\) utiliza el producto euclídeo de la coordenatización equilibrada del estado inicial; un cambio de coordenatización transporta también ese producto métrico.

Un transporte \(H\) ortogonal y simpléctico puede tener \(D\neq0\) y, a la vez, conservar el producto gaussiano isotrópico, con \(\nu=1\). El complemento operatorio y el entrelazamiento del estado tienen, por tanto, criterios distintos. En esta preparación, la deformación relativa no isométrica caracteriza exactamente el caso \(\nu>1\).

\subsubsection{Evaluación exacta del lazo áureo}

Para \eqref{md:antmem:cuantica:eq:5.2},

\begin{equation}
\operatorname{tr}(H_nH_n^{\mathsf T})-2
=n^2(2n^2-2n+5),\qquad
\nu_n^2=1+\frac8{81}n^2(2n^2-2n+5).
\label{md:antmem:cuantica:eq:5.8}
\end{equation}

La identidad resulta de sumar los cuadrados de las cuatro entradas. Para \(n=1\), la suma es siete y

\begin{equation}
\boxed{\nu_1=\frac{11}{9},\quad
\operatorname{Tr}\rho_{\rm vis}^2=\frac9{11},\quad
\bar n_{\rm W}=\frac19,\quad
\frac{S_{\rm vis}}{k_B}=\frac{10}{9}\log10-\log9.}
\label{md:antmem:cuantica:eq:5.9}
\end{equation}

El teorema~\ref{md:antmem:gauss:thm:espectro-gaussiano} demuestra que el estado gaussiano centrado de un modo con valor simpléctico \(\nu\) es unitariamente equivalente a la densidad geométrica de \eqref{md:antmem:gauss:eq:estado-geometrico}, cuya ocupación media es \(\bar n_{\rm W}=(\nu-1)/2\). Para \(\nu=11/9\), sus autovalores son

\begin{equation}
\lambda_j=\frac9{10}\,10^{-j},\qquad j=0,1,2,\ldots.
\label{md:antmem:cuantica:eq:5.10}
\end{equation}

La serie suma uno. Su cuadrado suma
\((9/10)^2/(1-10^{-2})=9/11\). Finalmente,

\[
-\sum_j\lambda_j\log\lambda_j
=-\log(9/10)+\log10\sum_jj\lambda_j
=\frac{10}{9}\log10-\log9.
\]

La reducción gaussiana y su espectro están demostrados en la sección~\ref{md:antmem:gauss:gaussianas-espectro}; las dos series anteriores evalúan exactamente la pureza y la entropía de esta realización. Como el estado conjunto es puro, \(S_{\rm vis}/k_B\) es también su entropía adimensional de entrelazamiento. La transformación global es unitaria, mientras la restricción al primer modo tiene entropía positiva.

El resultado corresponde a la preparación y al transporte especificados. Su identificación con una medición del vacío requiere la realización experimental de ese protocolo. La entropía de una reducción y el calor transferido se mantienen como magnitudes distintas.
